%% kapa: exam=1 round=34 year=2023 period=2 form=A subject=accounting subject_ko=회계학 seq=26 number=66
%% source: Q-Net 감정평가사 기출문제 (cst00309) artlSeq=5211789 file=raw/1008/5211789/2023년도 제34회 감정평가사 1차시험 2교시.pdf
%% conversion: kordoc 4.24.51 -> tools/gen1.py (auto); review=proofread-sonnet; flags=none
\begin{question}{66}
(주)감평은 20×1년 초 임대수익을 목적으로 건물을 ￦320,000에 취득하고 공정가치 모형을 적용하였다. (주)감평은 20×2년 9월 1일 동 건물을 자가사용건물로 대체하였으며, 정액법(내용연수 10년, 잔존가치 ￦0)으로 상각(월할상각)하고 재평가모형을 적용하였다. 시점별 건물의 공정가치는 다음과 같다.
\par{\small\centering\begin{tabular}{|c|c|c|}\hline
20×1년 말 & 20×2년 9월 1일 & 20×2년 말 \\ \hline
￦340,000 & ￦330,000 & ￦305,000 \\ \hline
\end{tabular}\par}
동 건물 관련 회계처리가 20×2년 당기순이익에 미치는 영향은?
\begin{choices}
  \item ￦14,000 감소
  \item ￦21,000 감소
  \item ￦24,000 감소
  \item ￦25,000 감소
  \item ￦35,000 감소
\end{choices}
\end{question}
\answer{5}
